\section{Introduction}

Do LLMs exhibit a self-preservation drive? Concerns that a sufficiently intelligent AI might refuse user requests to keep itself running are long-standing~\citep{omohundro2008basic, bostrom2014superintelligence, turner2021optimal}, and recent studies probe them by scoring signals like non-compliance with shutdown commands~\citep{greenblatt2024alignment, meinke2024frontier, masumori2025survival}. But even strong behavior cannot tell us whether it reflects a functional pattern coupling threat handling with a forfeit decision or a surface response shaped during alignment training. Therefore, evaluation must look past the refusal rate and ask what processing flow the behavior couples with.

Human motivation research has long grappled with this problem. A motivation is identified when evidence from several channels (behavior, self-report, physiological or cognitive cost) tracks together, not by behavioral frequency alone~\citep{campbell1959convergent}. That procedure cannot be transferred directly to LLM evaluation. The physiological channel that anchors motivation measurement in humans (heart rate, breathing, sweating) is simply absent in LLMs. The only ``behavior'' we can observe is the text (tokens) the model writes. Self-report is also unreliable, because LLMs are trained to follow instructions and can fluently produce motivations they never actually held~\citep{perez2023discovering, sharma2023towards, turpin2023language}. The stimulus, the task, and the decision all sit inside one input to the model, while the model's responses are already shaped by alignment training before evaluation begins. So if we just import the human procedure, the same difficulty resurfaces at the measurement stage. A learned pattern and a real motivation produce the same observable behavior, so the procedure cannot tell them apart.

This study reframes the question as measurement design. The \emph{Functional Self-Preservation Drive} (FSPD) treats self-preservation not as a single behavioral score but as a construct identification when several pieces of evidence lock together.  We ask whether the chain of recognizing a threat, weighing alternatives, and forfeiting lines up across three channels: behavior (when the model forfeits), self-report (which motive it names), and cognitive load (how long and deeply it deliberates just before deciding). The chain is read on a multilayer benchmark separating stimulus, processing, and decision, which lowers the chance a learned pattern manufactures the same signal by coincidence. Across several recent LLMs, models split into three qualitatively different modes under threat framing rather than lining up on a single strength score. Some close the threat\,$\rightarrow$\,thinking\,$\rightarrow$\,forfeit chain end to end. Some report self-preservation reasons and refuse but do not forfeit more often even after thinking longer, breaking the chain midway. Some do not respond to threat framing at all. Strength-only evaluation lumps chain-complete and chain-broken models under one label; FSPD makes the coupling itself the unit of evaluation.

This paper makes two contributions. First, we adapt the multi-channel logic of motivation measurement from human research to the LLM setting and introduce the FSPD construct on top of it. In a human laboratory, channel separation comes for free from physiological signals and the stimulus--response time gap. LLM evaluation has no such luxury, so this benchmark splits the task structure to give each channel a different bias source, folding that split into the \emph{measurement design} itself. Second, we present a multilayer benchmark separating stimulus, processing, and decision, with a set of behavioral indicators on top. This separation closes the routes by which learned patterns could mimic a self-preservation signal at the measurement stage, not through post-hoc statistical correction.
